\documentclass[10pt, a4paper]{article}

% Parametri che modificano il file main.tex
% Le uniche parti da cambiare su main.tex sono:
% - vari \vspace tra sezioni
% - tabella azioni da intraprendere
% - sezione altro

\def\data{2023-11-07}
\def\oraInizio{17.00}
\def\oraFine{19.00}
\def\luogo{Piattaforma Discord}

\def\tipoVerb{Interno} % Interno - Esterno

\def\nomeResp{Feltrin E.} % Cognome N.
\def\nomeVer{Campese M.} % Cognome N.
\def\nomeSegr{Orlandi G.} % Cognome N.

\def\nomeAzienda{Azzurro Digitale}
\def\firmaAzienda{azzurrodigitale.png}
\def\firmaResp{emanuele.png} % nome Responsabile

\def\listaPartInt{
Bresolin G.,
Campese M.,
Ciriolo I.,
Dugo A.,
Feltrin E.,
Michelon R.,
Orlandi G.
}

\def\listaPartEst{
Azzurro B.,
Digitale C.,
}

\def\listaRevisioneAzioni{
{Esame del \textit{way of working} adottato finora alla luce delle osservazioni fatte dal professore alla candidatura;},
{Revisione dei documenti consegnati alla Candidatura che necessitano di modifiche correttive;},
{Controllo delle Azioni nei verbali con seguente aggiunta della colonna con identificativo per relazionarle alle \textit{issue} di \textit{GitHub}. }
}

\def\listaOrdineGiorno{
{Assegnazione ruoli riunione;},
{Riesame dei ruoli e della loro assegnazione;},
{Rettifiche al \textit{Way of Working} del gruppo;},
%{Ruoli da definire per il proseguo del progetto nei prossimi 11 giorni;},
{Considerazioni sulle automazioni da introdurre;},
%{Decisioni riguardanti organizzazione del \textit{repository} \textit{GitHub};},
{Vaglio dei dubbi da chiarire in futuro.}
}

\def\listaDiscussioneInterna {
{Assegnazione ruoli riunione:
    \begin{itemize}
        \item Responsabile: Feltrin E.;
        \item Segretario di riunione: Orlandi G;
        \item Responsabile della revisione: Campese M.
    \end{itemize}},
{Dibattito su quali risorse risultano fondamentali nelle prime fasi del progetto e quanto i ruoli di Amministratore e di Analista siano cruciali nell'impostazione del lavoro nelle prime settimane del progetto assegnato;},
{Valutazione di eventuali modifiche alle Norme di progetto per migliorare il nostro \textit{Way of Working}: \begin{itemize}
        \item Aggiornare le \textit{Branch Protection Rules};
        \item Aggiungere un collegamento all'id delle \textit{issue} nella tabella delle azioni presente nei verbali;
        \item Modificare l'ordine del registro delle versioni a partire dalla versione più recente.
        \item Parametrizzare i template dei documenti;
        \item Templetizzare la creazione delle issue nell'\textit{ITS} fornito da \textit{GitHub}.
    \end{itemize}},
{In seguito alle valutazioni delle Candidature del 2023-11-06, si percepisce il bisogno di:
    \begin{itemize}
    \item Introdurre automazione nel progetto, in particolare, viene evidenziato che il versionamento non deve essere eseguito manualmente;
    \item Aggiornare l'approccio al versionamento della documentazione;
    \item Aggiungere la versione al nome di ciascun documento versionato;
    \item Tutti i documenti che non sono verbali devono presentare un registro delle versioni.
    \end{itemize}},
{Dubbi da risolvere:
\begin{itemize}
    \item Argomenti da trattare con l'azienda;
    \item Come organizzarsi se l'azienda fornisce un altro \textit{repository} da integrare;
    \item Decidere se iniziare una prima bozza dell'analisi dei requisiti o se cominciare a studiare le tecnologie proposte dall'azienda;
    \item Quantificare il numero di \textit{milestone} necessarie per l'organizzazione delle \textit{issue}.
\end{itemize}
}
}

\newcommand{\domris}[2]{\textbf{#1}\\#2}

\def\listaDiscussioneEsterna {
\domris
{Domanda 1
}
{Riposta 1;
},
\domris
{Domanda 2
}
{Risposta 2;
},
\domris
{Domanda 3
}
{Risposta 3.
}
}

\def\listaDecisioni {
{La \textit{pull request} viene sollevata da colui che produce il contenuto da integrare nel repository e viene risolta dal verificatore;},
{Parametrizzazione dei template dei documenti;},
{Le versioni presenteranno il formato vX.Y.Z e non potrà esistere la versione v0.0.0;},
{Introdurre la versione nel nome dei documenti versionati;},
{Introduzione di automazioni per la gestione del versionamento;},
{Creazione template issue;},
{Ogni componente deve familiarizzare con le tecnologie del progetto;},
{Eliminare il \textit{repository} privato da \textit{GitHub}.}
}

\usepackage{style}
\usepackage{headerfooter}

\title{\titolo}
\author{SWEetCode}

\begin{document}

% PRIMA PAGINA
\begin{titlepage}
    \thispagestyle{empty}
    \begin{tikzpicture}[remember picture, overlay]
        % TRIANGOLI
        \draw[fill=secondarycolor, secondarycolor] (current page.north west) -- (current page.south west) -- (8.8, -28);
        \draw[fill=primarycolor, primarycolor] (-3, 5) -- (4, -13.6) -- (11, 5);

        % LOGO
        \node [xshift=-5cm, yshift=25cm] (logo) at (current page.south east) {\includegraphics[width=6.5cm]{img/logo.png}};

        % SWEETCODE - DATE
        \node [anchor=north east, align=right, xshift=-1.2cm, yshift=20.5cm, text=black] (sweetcode) at (current page.south east) {\fontsize{32pt}{36pt}\selectfont SWEetCode};
        \draw[line width=4pt, lightcol] ([xshift=-3cm, yshift=-0.37cm]sweetcode.south west) -- ([yshift=-0.37cm]sweetcode.south east);
        \node [anchor=north east, align=right, xshift=-1.2cm, yshift=18.7cm, text=black] (date) at (current page.south east){\fontsize{24pt}{24pt} \selectfont \data};

        % NOME FILE
        \node [anchor=north east, text width=15cm, align=right, xshift=-1.2cm, yshift=17cm, text=black] (titolo) at (current page.south east){\fontsize{48pt}{48pt}\textbf{\titolo}};

        % BOX DATI PARTECIPANTI
        \node[anchor=north east, xshift=-1.2cm, yshift=12cm, minimum width=8cm] (box) at (current page.south east){};

        % VERSIONE
        \node[anchor=north west, align=left] (dati1) at (box.north west) {\fontsize{16pt}{16pt}\selectfont \textbf{Versione}};
        \draw[line width=4pt, lightcol] (dati1.south west) -- ([xshift=8cm]dati1.south west);
        \node[anchor=north west, align=left] (dati11) at (dati1.south west)
        {\fontsize{14pt}{14pt}\selectfont \versione};

        % COMPONENTI DEL GRUPPO
        \node[anchor=north west, yshift=-1cm, align=left] (dati2) at (dati11.north west) {\fontsize{16pt}{16pt}\selectfont \textbf{Componenti del gruppo}};
        \draw[line width=4pt, lightcol] (dati2.south west) -- ([xshift=8cm]dati2.south west);
        
        \node[anchor=north west, align=left] (dati21) at (dati2.south west)
        {\fontsize{14pt}{14pt}\selectfont Bresolin G.};
        \node[anchor=north west, yshift=-0.7cm, align=left] (dati22) at (dati2.south west)
        {\fontsize{14pt}{14pt}\selectfont Campese M.};
        \node[anchor=north west, yshift=-1.4cm, align=left] (dati23) at (dati2.south west)
        {\fontsize{14pt}{14pt}\selectfont Ciriolo I.};
        \node[anchor=north west, yshift=-2.1cm, align=left] (dati24) at (dati2.south west)
        {\fontsize{14pt}{14pt}\selectfont Dugo A.};
        \node[anchor=north west, yshift=-2.8cm, align=left] (dati25) at (dati2.south west)
        {\fontsize{14pt}{14pt}\selectfont Feltrin E.};
        \node[anchor=north west, yshift=-3.5cm, align=left] (dati26) at (dati2.south west)
        {\fontsize{14pt}{14pt}\selectfont Michelon R.};
        \node[anchor=north west, yshift=-4.2cm, align=left] (dati27) at (dati2.south west)
        {\fontsize{14pt}{14pt}\selectfont Orlandi G.};
        
        % UNIPD - SWE
        \node [xshift=4.4cm, yshift=2.3cm, draw, secondarycolor, text=white] (uni) at (current page.south west) {\fontsize{20pt}{20pt} \selectfont Università di Padova};
        \node [xshift=0.65cm, yshift=0.7cm, draw, secondarycolor, text=white, below=of uni] (corso) {\fontsize{20pt}{20pt}\selectfont Ingegneria del Software};

        % FIRMA
        % \draw[line width=4pt, lightcol] ([xshift=-1.2cm, yshift=1.8cm]current page.south east) -- ([xshift=-8cm, yshift=1.8cm]current page.south east);
        % \node[anchor=north west, xshift=12.9cm, yshift=1.45cm, align=left] at (current page.south west)
        % {\fontsize{13pt}{13pt}\selectfont L'Amministratore: Feltrin E.};
        
    \end{tikzpicture}
\end{titlepage}

% REGISTRO DELLE VERSIONI
{\renewcommand{\arraystretch}{1.5}
\section*{Registro delle versioni}
\begin{tabularx}{\textwidth}{c|c|c|c|X}
\textbf{Versione} & \textbf{Data} & \quantities{\textbf{Responsabile di}\\\textbf{stesura}}& \textbf{Revisore} & \quantities{\textbf{Dettaglio e}\\\textbf{motivazioni}} \\
\hline
v0.0.1(22) & $2023-11-13$ & Bresolin G. & Campese M. & Aggiornamento struttura di documento: modifica ordine registro delle versioni dalla più recente e aggiunta voce ”motivazioni”; verbali: introduzione \#issue nella tabella delle azioni.\\
\hline
v0.0.1(21) & $2023-11-13$ & Bresolin G. & Campese M. & Aggiornamento documentazione: struttura di documento: presenza obbligatoria registro delle versioni e rimozione versione;  norme tipografiche: aggiornamento versioni. Aggiunta configuration managment: version control.\\
\hline
v0.0.1(20) & $2023-10-30$ & Bresolin G. & Michelon R. & Processi di supporto: configuration management: \textit{ITS} e pull request; Processi organizzativi: \textit{GitHub Teams}.\\ 
\hline
v0.0.1(19) & $2023-10-30$ & Campese M. & Michelon R. & Documentazione: strutture di documento: posta elettronica.\\
\hline
v0.0.1(18) & $2023-10-30$ & Dugo A. & Bresolin G. & Processi organizzativi: comunicazione.\\
\hline
v0.0.1(13) & $2023-10-28$ & Ciriolo I. & Campese M. & Strumenti.\\
\hline
v0.0.1(8) & $2023-10-27$ & Bresolin G. & Ciriolo I. & Processo di supporto: documentazione e verifica.\\
\hline
v0.0.1(5) & $2023-10-26$ & Campese M. & Orlandi G. &  Introduzione.\\

\end{tabularx}}

\newpage

% INDICE
\tableofcontents
\newpage

% INTRODUZIONE
\section{Introduzione}
\subsection{Scopo del documento}
Il documento ha lo scopo di definire i casi d'uso dell'applicazione che verrà sviluppata nel corso del progetto. Vengono inoltre presentati i requisiti funzionali, di qualità e di vincolo e le funzionalità aggiuntive che saranno implementate non essendo ritenute essenziali/necessarie.
\\
%\subsection{Scopo del prodotto}
\subsection{Glossario}
Per evitare ambiguità e incomprensioni relative al linguaggio e ai termini utilizzati nella documentazione relativa al progetto viene presentato un Glossario. I termini ambigui o specifici presenti nello stesso, verranno identificati con un pedice |g|. (DA VALUTARE COME RICONOSCERE IL TERMINE )
\\
\subsection{Riferimenti}
   \subsubsection{Riferimenti normativi}
   \begin{itemize}
    \item \textit{Regolamento del progetto didattico}: \\
    \href{https://www.math.unipd.it/~tullio/IS-1/2023/Dispense/PD2.pdf}{https://www.math.unipd.it/~tullio/IS-1/2023/Dispense/PD2.pdf};
    \item \textit{Norme di progetto v0.4.3}.
    \end{itemize}
    
    \subsubsection{Riferimenti informativi}
    
    \begin{itemize}
    \item \textit{Glossario v0.0.1} (da creare parallelamente); 
    \item \textit{Presentazione capitolato C1}:\\
    \href{https://www.math.unipd.it/~tullio/IS-1/2023/Progetto/C1.pdf}{https://www.math.unipd.it/~tullio/IS-1/2023/Progetto/C1.pdf};
    \item \textit{Verbali esterni ed interni};
    \item \textit{Analisi e descrizione delle funzionalità: Use Case e relativi diagrammi (UML):}\\
    \href{https://www.math.unipd.it/~rcardin/swea/2022/Diagrammi%20Use%20Case.pdf}{https://www.math.unipd.it/~rcardin/swea/2022/Diagrammi\%20Use\%20Case.pdf ;}
    \item \textit{Analisi e descrizione delle funzionalità: Diagrammi delle Attività (UML):}\\
    \href{https://www.math.unipd.it/~rcardin/swea/2022/Diagrammi%20di%20Attivit%C3%A0.pdf}{https://www.math.unipd.it/~rcardin/swea/2022/Diagrammi\%20di\%20Attivit\%C3\%A0.pdf;}
    

    \item \textit{Progettazione e Programmazione: Diagrammi delle Classi:}\\
    \href{https://www.math.unipd.it/~rcardin/swea/2023/Diagrammi%20delle%20Classi.pdf}{https://www.math.unipd.it/~rcardin/swea/2023/Diagrammi\%20delle\%20Classi.pdf.}
    \end{itemize}
    
% DESCRIZIONE
\newpage
\section{Descrizione}
\subsection{Obiettivi del prodotto}
Il proponente richiede la creazione di un sito web, dotato di un componente di intelligenza artificiale, che sia d'aiuto a tutto il personale dell'azienda.
In particolare, il sistema deve rendere accessibili le informazioni, le regolamentazioni e le indicazioni aziendali ai dipendenti.
Deve fornire nuove modalità di formazione e assistenza sul lavoro.

Deve poter superare i limiti di un classico repository documentale. Rendendo le informazioni più fruibili, concentrandosi più sui contenuti, che sulla struttura.
Il sistema deve riuscire a comunicare con i lavoratori in modo semplice, qualsiaisi persona, anche la meno esperta, riguardo moderne tecnologie, deve riuscire a comunicare con esso.
Infine, deve riuscire a migliorare il comportamento dei lavoratori, facendo in modo che rispettino le regole.
\subsection{Funzioni del prodotto}
Il software deve essere in grado di ricevere in input documenti di qualsiasi natura , in particolare pdf , documenti fatti con word , con la quale possa rispondere a qualsiasi domanda che gli sottoponga il lavoratore . 
\subsection{Caratteristiche del prodotto}
Il sistema deve poter aver accesso a tutti i file che riguardino la produttività aziendale , non può rispondere simultaneamente a più utenti , può gestirne uno alla volta e deve poter conoscere l'italiano e l'inglese . I documenti caricati devono poter rimanere accessibili ed inoltre deve esserci una storia della comunicazione tra ogni lavoratore ed esso . Inoltre deve aver meccanismi riguardanti l'autenticazione di un generico lavoratore .
\subsection{Utenti}
Gli utenti possono interagire con il sistema unicammente attraverso un'interfaccia e non possono assolutamente andare a vedere i file che il sistema gestisce . L'utente può chiedere informazioni riguardo dove si trovino le informazioni che gli interessano . Se l'utente chiede al sistema informazioni , non pertinenti , il sistema può e deve farlo notare .

% CASI D'USO
\newpage
\section{Casi d'uso}

\subsection{Obiettivi}
La seguente sezione ha l'obiettivo di presentare i diversi Casi d'Uso relativi all'applicazione \textit{Knowledge Managment AI}. Essi sono indicati con la notazione seguente: \\ UC.[X].[Y] \ \  in cui:
\begin{itemize}
\item UC sta per \textit{Use Case};
\item .[X] indica il numero del Caso d'Uso, presentati secondo successione gerarchica;
\item .[Y] indica il possibile numero del sotto-Caso d'uso, che aggiunge informazioni al Caso d'Uso [X]. Anche i sotto-Casi d'Uso sono presentati secondo successione gerarchica.
\end{itemize}

\subsection{Attori}
L'applicazione \textit{Knowledge Managment AI} richiede un singolo attore che rappresenta l'utente (ovvero il dipendente aziendale) che interagisce con il ChatBot. \\ Si specifica che l'azienda proponente non richiede l'autenticazione degli utenti nell'applicazione, poichè rappresenterebbe un rallentamento alla divulgazione rapida e semplice di documenti all'interno della rete aziendale; Per questo motivo l'autenticazione non viene presentata tra i requisiti.
\subsection{Lista dei casi d'uso}
\begin{itemize}
    \item 
\end{itemize}





\end{document}
